% LỜI TỰA --- Quyển XX
\chapter*{Lời Tựa}
\addcontentsline{toc}{chapter}{Lời Tựa}
\markboth{Lời Tựa}{Lời Tựa}

\begin{truyen}{Lời Tựa}{50 Sentences: Funny but Profound}
\chuhoa{G}{iáo viên bắt đầu bộ sách này với một câu hỏi đơn giản: làm thế nào để mỗi buổi học có thể để lại thứ gì đó đáng nhớ hơn là bài thi?}
\textit{(A teacher began this series with a simple question: how can each class leave behind something more memorable than a test?)}

Những câu nói hài hước của thầy giáo trong lớp đôi khi khiến cả lớp cười~--- nhưng cười xong, học sinh nhớ.
\textit{(Humorous teacher sayings in class sometimes make everyone laugh~--- but after the laugh, students remember.)}

Quyển hai mươi gom 50 câu vừa gây cười vừa gợi suy nghĩ: ngắn, dễ nhớ, và sau tiếng cười vẫn để lại bài học. Mong thầy cô thêm vài câu vui mà sâu; học trò thêm vài kỷ niệm vừa cười vừa nghĩ.
\textit{(Volume twenty gathers 50 sentences that both amuse and make you think: short, memorable, and after the laugh they still leave a lesson. May teachers add a few lines that are funny and deep; may students keep a few memories that mix laughter and thought.)}

Bộ sách <<Truyện Tích Cảm Hứng Song Ngữ>> gồm 24 quyển, mỗi quyển một chủ đề riêng. Quyển XX này mang chủ đề: \textbf{thầy cười mà học trò nghĩ~--- hài hước đúng cách}.
\textit{(The series ``Inspiring Bilingual Tales'' contains 24 volumes, each with its own theme. Volume XX focuses on: \textbf{thầy cười mà học trò nghĩ~--- hài hước đúng cách}.})
\end{truyen}

\begin{baihoc}
Hài hước đúng cách không làm mất uy tín~--- nó làm học trò gần thầy và nhớ lâu hơn.
\textit{(The right kind of humor doesn't undermine authority~--- it brings students closer and makes them remember longer.)}
\end{baihoc}

\begin{ghinhoanh}
\textbf{Short English to remember:}

Laughter that leaves a lesson is the best kind.

Humor and wisdom are not opposites.
\end{ghinhoanh}

\begin{tuhocnhanh}
1. Câu nói vừa buồn cười vừa sâu sắc nào em vẫn nhớ sau khi đọc xong? \textit{(Which funny yet deep sentence do you still remember after finishing the book?)}

2. Em có thầy cô nào hay dùng hài hước để dạy không? Kể một ví dụ? \textit{(Do you have a teacher who uses humor to teach? Give one example.)}

3. Tự em nghĩ được câu hài hước mà vẫn có bài học không? Thử viết một câu. \textit{(Can you write a funny sentence that still carries a lesson? Try one.)}
\end{tuhocnhanh}

\begin{center}
\textit{\color{teal}``Cười xong nhớ mới là hay.''}
\end{center}

\begin{flushright}
\textbf{Tôi Nguyễn Văn Sang}\\
\textit{GV THPT Nguyễn Hữu Cảnh - TP HCM}
\end{flushright}

\ngancach
